% !TEX root = ../main.tex
\section{Hạn chế và hướng phát triển}\label{sec:limit}

\paragraph{Hạn chế.}
\begin{enumerate}[nosep]
  \item Chưa có nhãn chữ thật, nên CRNN chưa được fine-tune trên biển thật và chưa có số liệu ablation end-to-end.
  \item Dữ liệu tổng hợp dùng font Hershey, khác font biển số Việt Nam.
  \item Chưa tối ưu cho ảnh ban đêm hay camera hồng ngoại.
  \item Hướng vào/ra hiện cố định theo camera.
\end{enumerate}

\paragraph{Hướng phát triển, theo thứ tự ưu tiên.}
\begin{enumerate}[nosep]
  \item \textbf{Gán nhãn chữ bán tự động.} PaddleOCR đọc trước trên các khung biển có sẵn trong nhãn YOLO, sau đó người kiểm tra và sửa. Từ đó tạo \code{texts.csv} và \code{test\_texts.csv}, rồi fine-tune CRNN (bước 4b) và chạy ablation.
  \item Sinh dữ liệu tổng hợp bằng font gần giống font biển số Việt Nam, để giảm khoảng cách miền ngay từ giai đoạn huấn luyện trước.
  \item Bổ sung các cặp nhầm đặc thù của font Việt Nam (ví dụ \code{H}$\to$\code{4}) vào hậu xử lý, dựa trên phân tích lỗi thực tế.
  \item Xuất ONNX với kích thước đầu vào động để tăng tốc trên CPU. Thử TensorRT FP16 khi triển khai trên GPU.
  \item Suy ra hướng vào/ra từ quỹ đạo theo dõi thay vì cố định theo camera.
\end{enumerate}
